\section{Optimality}\label{sec:optimality}

The bound of \cref{thm:main} is best possible up to the factors it hides. Driemel et
al.~\cite{DMOSW25} prove that every randomized algorithm that estimates $\wt(\MST(P))$ to within a
constant multiplicative factor, with probability at least $2/3$, must make $\Omega(n^{1/3})$
orthogonal range-counting queries; their lower bound is built from a family of planar instances that a
range-counting algorithm cannot distinguish without $\Omega(n^{1/3})$ queries, and it applies a fortiori
to $(1\pm\eps)$-estimation. Combining this with \cref{thm:main} gives the matching bound.

\begin{corollary}\label{cor:tight}
The randomized orthogonal range-counting query complexity of estimating the planar Euclidean MST weight to
a $(1\pm\eps)$ factor, for any fixed $\eps\in(0,1)$ and inputs $P\subseteq[\Delta]^2$ with $\Delta=O(n)$, is
$n^{1/3}$ up to polylogarithmic and $\eps$-dependent factors.
\end{corollary}

\begin{remark}[Model agreement]\label{rem:model}
The match is between two bounds in the same model, so it is worth stating the model exactly. In both
\cref{thm:main} and the lower bound of~\cite{DMOSW25}, the algorithm is randomized, knows $n$ and
$\Delta$, and adaptively issues axis-aligned integer rectangles $R\subseteq[\Delta]^2$, each answered by
the exact count $\abs{P\cap R}$; the cost is the number of such queries and there is no other access to
$P$. Our algorithm uses only queries of this form (every rectangle below is a union of grid cells with
integer corners), and the lower bound holds against all such adaptive randomized algorithms with success
probability $\ge 2/3$. The two bounds therefore compose, and \cref{cor:tight} is a statement about this one
oracle, not an artifact of differing query primitives or success criteria.
\end{remark}

The upper and lower bounds meet at the same exponent and under the same oracle, so the role of range
counting in Euclidean MST-weight estimation is fixed: the $\Otil(\sqrt n)$ of~\cite{DMOSW25} was not the
true rate, and no further polynomial improvement in the query exponent is available. The remaining slack
is the polylogarithmic and $\eps$-dependent overhead carried by the estimator of \cref{sec:support} and
the death-time sampler of \cref{sec:dt}.

\paragraph{Why $n^{1/3}$.}
The exact lower-bound construction is that of~\cite{DMOSW25}, which we take as given; here we only record
the conceptual reading that our balance point makes transparent, namely why the exponent is $1/3$ and not
$1/2$. Estimating the MST weight to a constant factor amounts to resolving the cluster structure at the
scale where it carries a constant fraction of the weight. At a grid scale with $K$ occupied cells, the
weight lives in the number of \emph{components} the cells form, and a single range count returns a
population, not a component count: it cannot, on its own, tell whether the points in a box form one
cluster or many. The number of scales and cells at which the answer is genuinely in doubt is governed by
the support size, which the regularization fixes at $K_0=\Theta(n^{2/3})$, and resolving a structure with
$\Theta(n^{2/3})$ independently-uncertain pieces against an oracle that only counts has a
$\sqrt{\,\cdot\,}$-type sampling cost in the standard query-complexity sense~\cite{AroraBarak09},
$\Theta(\sqrt{n^{2/3}})=\Theta(n^{1/3})$. This is the same $\sqrt K$
that the leader estimator pays on the upper-bound side, and the coincidence is not accidental: $1/3=\tfrac12
\cdot\tfrac23$ is the sampling exponent against an $n^{2/3}$-sized uncertain support.

Our algorithm spends its budget at exactly this granularity. It regularizes to a support of
$K_0=\Theta(n^{2/3})$ cells, where the support estimator's two terms $\sqrt K$ and $n/K$ are equal, and
reads the cluster structure there in $\Otil(n^{1/3})$ queries. The spatial leader estimator is what lets
the $\sqrt K$ term appear in place of a $K$ term; without it the same reduction would pay one query per
occupied cell and reproduce the $\Otil(\sqrt n)$ bound. In this heuristic sense the leader estimator pays
the $\sqrt{\,\cdot\,}$ sampling rate against the $n^{2/3}$-sized support, the same exponent the lower bound
extracts; the formal statement is the matching of \cref{thm:main} and the $\Omega(n^{1/3})$ bound
of~\cite{DMOSW25}, which is \cref{cor:tight}.
